% ============================================================
%  Lecture Notes Template — Monotone Classification
%  Based on: "Monotone Classification with Relative Approximations"
%  Usage: Copy any box below and paste wherever you need it.
% ============================================================

\documentclass[11pt, a4paper]{article}

% ---------- packages ----------
\usepackage[margin=1in]{geometry}
\usepackage{amsmath, amssymb, amsthm}
\usepackage{mathtools}
\usepackage[T1]{fontenc}
\usepackage{lmodern}
\usepackage{microtype}
\usepackage{hyperref}
\usepackage{xcolor}
\usepackage{mdframed}          % coloured boxes
\usepackage{enumitem}
\usepackage{thmtools}
\usepackage{graphicx}
\usepackage{booktabs}          % nice tables

% ---------- colours (simple, soft palette) ----------
\definecolor{defblue}{RGB}{220, 232, 245}   % light steel blue
\definecolor{thmgreen}{RGB}{220, 240, 225}  % light sage green
\definecolor{lemyellow}{RGB}{255, 248, 220} % light cream
\definecolor{propgray}{RGB}{235, 235, 235}  % light gray
\definecolor{noteorange}{RGB}{255, 237, 213}% light peach
\definecolor{rulecol}{RGB}{100, 130, 160}   % muted blue — used for box borders

% ---------- box styles (mdframed) ----------
\mdfdefinestyle{defstyle}{
  backgroundcolor = defblue,
  linecolor       = rulecol,
  linewidth       = 1pt,
  innertopmargin  = 6pt,
  innerbottommargin = 6pt,
  innerleftmargin = 8pt,
  innerrightmargin = 8pt,
  skipabove       = 8pt,
  skipbelow       = 8pt,
}

\mdfdefinestyle{thmstyle}{
  backgroundcolor = thmgreen,
  linecolor       = rulecol,
  linewidth       = 1pt,
  innertopmargin  = 6pt,
  innerbottommargin = 6pt,
  innerleftmargin = 8pt,
  innerrightmargin = 8pt,
  skipabove       = 8pt,
  skipbelow       = 8pt,
}

\mdfdefinestyle{lemstyle}{
  backgroundcolor = lemyellow,
  linecolor       = rulecol,
  linewidth       = 1pt,
  innertopmargin  = 6pt,
  innerbottommargin = 6pt,
  innerleftmargin = 8pt,
  innerrightmargin = 8pt,
  skipabove       = 8pt,
  skipbelow       = 8pt,
}

\mdfdefinestyle{propstyle}{
  backgroundcolor = propgray,
  linecolor       = rulecol,
  linewidth       = 1pt,
  innertopmargin  = 6pt,
  innerbottommargin = 6pt,
  innerleftmargin = 8pt,
  innerrightmargin = 8pt,
  skipabove       = 8pt,
  skipbelow       = 8pt,
}

\mdfdefinestyle{notestyle}{
  backgroundcolor = noteorange,
  linecolor       = rulecol,
  linewidth       = 1pt,
  innertopmargin  = 6pt,
  innerbottommargin = 6pt,
  innerleftmargin = 8pt,
  innerrightmargin = 8pt,
  skipabove       = 8pt,
  skipbelow       = 8pt,
}

% ---------- theorem-like environments ----------
\newtheoremstyle{boxed}{}{}{\normalfont}{}{\bfseries}{.}{.5em}{}

\theoremstyle{boxed}
\newtheorem{innerdefinition}{Definition}[section]
\newtheorem{innertheorem}{Theorem}[section]
\newtheorem{innerlemma}{Lemma}[section]
\newtheorem{innerproposition}{Proposition}[section]
\newtheorem{innercorollary}{Corollary}[section]
\newtheorem{innerproblem}{Problem}[section]
\newtheorem{innerexample}{Example}[section]
\newtheorem{innerremark}{Remark}[section]

% Wrap each in its coloured box
% ── DEFINITION ──────────────────────────────────────────────
\newenvironment{definition}
  {\begin{mdframed}[style=defstyle]\begin{innerdefinition}}
  {\end{innerdefinition}\end{mdframed}}

% ── THEOREM ─────────────────────────────────────────────────
\newenvironment{theorem}
  {\begin{mdframed}[style=thmstyle]\begin{innertheorem}}
  {\end{innertheorem}\end{mdframed}}

% ── LEMMA ───────────────────────────────────────────────────
\newenvironment{lemma}
  {\begin{mdframed}[style=lemstyle]\begin{innerlemma}}
  {\end{innerlemma}\end{mdframed}}

% ── PROPOSITION ─────────────────────────────────────────────
\newenvironment{proposition}
  {\begin{mdframed}[style=propstyle]\begin{innerproposition}}
  {\end{innerproposition}\end{mdframed}}

% ── COROLLARY ───────────────────────────────────────────────
\newenvironment{corollary}
  {\begin{mdframed}[style=thmstyle]\begin{innercorollary}}
  {\end{innercorollary}\end{mdframed}}

% ── PROBLEM ─────────────────────────────────────────────────
\newenvironment{problem}
  {\begin{mdframed}[style=defstyle]\begin{innerproblem}}
  {\end{innerproblem}\end{mdframed}}

% ── EXAMPLE ─────────────────────────────────────────────────
\newenvironment{example}
  {\begin{mdframed}[style=notestyle]\begin{innerexample}}
  {\end{innerexample}\end{mdframed}}

% ── REMARK ──────────────────────────────────────────────────
\newenvironment{remark}
  {\begin{mdframed}[style=notestyle]\begin{innerremark}}
  {\end{innerremark}\end{mdframed}}

% Plain proof (no box needed — standard amsthm)
% \begin{proof} ... \end{proof}  works out of the box.

% ---------- common math shorthands ----------
\newcommand{\R}{\mathbb{R}}
\newcommand{\N}{\mathbb{N}}
\newcommand{\Z}{\mathbb{Z}}
\newcommand{\Hmon}{\mathbb{H}_{\mathrm{mon}}}
\newcommand{\err}{\mathrm{err}}
\DeclareMathOperator{\E}{\mathbf{E}}
\renewcommand{\Pr}{\mathbf{Pr}}
\newcommand{\kstar}{k^{*}}
\newcommand{\dom}{\succ}        % dominates
\newcommand{\predom}{\prec}
\newcommand{\preceqdom}{\preceq}

% ============================================================
\begin{document}

% ---------- TITLE BLOCK ----------
\title{\textbf{Monotone Classification with Relative Approximations}\\[4pt]
       \large Lecture Notes for SC2320 Data Analytics and Mining \\
       }
% \author{Your Name \\ \small Your University}

\date{\today}
\maketitle

% \begin{abstract}
%   Brief summary of what these notes cover.
%   Replace this with your own abstract or remove this block entirely.
% \end{abstract}

\section*{Instructions to readers}
The flow of this set of notes is in tandem to the flow of the original paper.
It will also go in the same flow as the presentation slides.
Concepts/Material/Information that is auxiliary to the core ideas of this paper will be outlined at the end of each section.

% ============================================================
\section*{Acronym Reference}
% ============================================================

\begin{table}[ht]
  \centering
  \begin{tabular}{ll}
    \toprule
    \textbf{Acronym} & \textbf{Full Form} \\
    \midrule
    ACR1 & Acronym One \\
    ACR2 & Acronym Two \\
    ACR3 & Acronym Three \\
    \bottomrule
  \end{tabular}
\end{table}

% \tableofcontents



\newpage

% ============================================================
% COPY-PASTE ENVIRONMENT TEMPLATES
% Uncomment any block and paste it wherever you need it.
% ============================================================

% % ──────────────────────────────────────────────
% %  BOX 1 — DEFINITION
% %  Copy from \begin{definition} to \end{definition}
% % ──────────────────────────────────────────────
% \begin{definition}[Monotone Classifier]
%   A classifier $h : \R^d \to \{-1,1\}$ is \emph{monotone} if
%   $h(p) \ge h(q)$ for every $p,q \in \R^d$ satisfying $p \dom q$.
% \end{definition}

% % ──────────────────────────────────────────────
% %  BOX 2 — THEOREM
% %  Copy from \begin{theorem} to \end{theorem}
% % ──────────────────────────────────────────────
% \begin{theorem}[Upper Bound]\label{thm:upper}
%   For any $\epsilon > 0$, there exists an algorithm that finds a
%   $(1+\epsilon)$-approximate monotone classifier using
%   $O\!\left(\frac{\log n}{\epsilon}\right)$ label queries.
% \end{theorem}

% \begin{proof}
%   Write your proof here. The $\square$ is inserted automatically.
% \end{proof}

% % ──────────────────────────────────────────────
% %  BOX 3 — LEMMA
% %  Copy from \begin{lemma} to \end{lemma}
% % ──────────────────────────────────────────────
% \begin{lemma}\label{lem:key}
%   For any input $P$ and monotone classifier $h \in \Hmon$,
%   \[
%     \err_P(h) \;\ge\; \kstar.
%   \]
% \end{lemma}

% \begin{proof}
%   Proof of the lemma.
% \end{proof}

% % ──────────────────────────────────────────────
% %  BOX 4 — PROPOSITION
% %  Copy from \begin{proposition} to \end{proposition}
% % ──────────────────────────────────────────────
% \begin{proposition}\label{prop:one}
%   If $P$ is monotone then $\kstar = 0$.
% \end{proposition}

% % ──────────────────────────────────────────────
% %  BOX 5 — COROLLARY
% %  Copy from \begin{corollary} to \end{corollary}
% % ──────────────────────────────────────────────
% \begin{corollary}\label{cor:one}
%   Theorem~\ref{thm:upper} implies that the algorithm of \ldots\ achieves
%   an error of at most $(1+\epsilon)\kstar$ with high probability.
% \end{corollary}

% % ──────────────────────────────────────────────
% %  BOX 6 — PROBLEM STATEMENT
% %  Copy from \begin{problem} to \end{problem}
% % ──────────────────────────────────────────────
% \begin{problem}[Monotone Classification with Relative Precision]\label{prob:main}
%   Given $\epsilon \ge 0$, find a monotone classifier whose error on $P$
%   is at most $(1+\epsilon)\kstar$.
%   The efficiency of an algorithm is measured by the number of elements probed.
% \end{problem}

% % ──────────────────────────────────────────────
% %  BOX 7 — EXAMPLE
% %  Copy from \begin{example} to \end{example}
% % ──────────────────────────────────────────────
% \begin{example}
%   Let $P = \{p_1, p_2, \ldots, p_n\} \subset \R^2$.
%   The optimal monotone classifier $h$ misclassifies only $p_1$, so $\kstar = 1$.
% \end{example}

% % ──────────────────────────────────────────────
% %  BOX 8 — REMARK / NOTE
% %  Copy from \begin{remark} to \end{remark}
% % ──────────────────────────────────────────────
% \begin{remark}
%   When $\epsilon = 0$ the problem reduces to finding an \emph{optimal}
%   monotone classifier, which is achievable in polynomial time after
%   revealing all labels.
% \end{remark}

% ============================================================
\section{Background and Motivation}
% ============================================================

\subsection{Entity Resolution}
Entity Resolution is a process in Data Preprocessing and Integration that aims to address the issue of matching and merging records corresponding to the same real-world object.

\begin{remark}
  Professor Tao's paper itself sometimes uses "Entity Matching" in the title and introduction to refer more broadly to the problem. For disambiguity in this set of notes, \textbf{"Entity Resolution"} will refer to the entire process while "Entity Matching" will refer to the process that the paper primarily focuses on.
\end{remark}
% \subsubsection{}



\subsection{The ER Pipeline}
% Overview of the pipeline stages:
%   Preprocessing → Blocking (mention LSH) → Comparison → Matching
% Flag Matching as the main focus of the paper.

\subsection{Matching in Depth}

\subsubsection{Why Human Labelling is Expensive}
% Discuss the cost of obtaining ground-truth labels from human annotators.

\subsubsection{The Binary Classification Objective}
% Frame matching as a binary classification problem: match (+1) vs non-match (−1).

\subsubsection{The Monotonicity Rule}
% Introduce the dominance relation and the monotonicity constraint on classifiers.

\subsubsection{Active Classification / Learning}
% Motivate active learning: query only the most informative labels to reduce cost.

\newpage

% ============================================================
\section{RPE Algorithm}
% ============================================================

\subsection{The Algorithm}
% Describe Random Probes with Elimination (RPE): iterative random probing
% and elimination of elements via the monotonicity constraint.

\subsection{Why It Works}
% Intuition and correctness argument behind the elimination strategy.

\subsection{Guarantees and Limitations}
% State the factor-of-2 approximation guarantee.
% Discuss where RPE falls short (e.g.\ dependence on width w).

\newpage

% ============================================================
\section{Coresets (RCC)}
% ============================================================

\subsection{Why RPE Isn't Enough}
% Explain the gap left by RPE and why a $(1+\varepsilon)$ guarantee requires more.

\subsection{The Coreset Technique}
% Introduce relative-comparison coresets (RCC): a weighted subset Z \subseteq P
% whose weighted error approximates err_P(h) to within relative error \varepsilon/4.

\subsection{Achieving the $1{+}\varepsilon$ Guarantee}
% Show how the coreset construction feeds into the full algorithm (Corollary 7)
% to obtain a $(1+\varepsilon)k^*$ error bound.

\newpage
% ============================================================
\section{Proofs and Lower Bounds}
% ============================================================

\subsection{Lower Bound Results}
% State Theorems 10, 13, 14: \Omega(n) for \varepsilon=0;
% lower bounds for constant and arbitrary approximation ratios.

\subsection{Near-Optimality of the Algorithms}
% Argue that the upper bounds from RPE and RCC match the lower bounds
% up to logarithmic factors, establishing near-optimality.

\subsection{Mathematical Machinery}
% Key tools used in proofs: concentration bounds, VC-dimension,
% disagreement coefficient, and the attrition-and-elimination game.

\newpage

% ============================================================
\section*{Summary}
% ============================================================
% Brief wrap-up: key takeaways, significance of results, and open questions.

% ============================================================
% \section*{References}
% ============================================================
\begin{thebibliography}{9}
  \bibitem{tao18} Y.~Tao, ``Monotone Classification,'' \emph{PODS'18}.
  \bibitem{tao21} Y.~Tao, ``Monotone Classification (Extended),'' \emph{PODS'21}.
\end{thebibliography}

\end{document}